\documentclass{article}
\usepackage[utf8]{inputenc}
\usepackage[T1]{fontenc}
\usepackage[margin=2cm]{geometry}
\usepackage{hyperref}
\sloppy
\begin{document}
\section*{Historical source}
Encyclopaedia Britannica scan, source row 2108. Parse 3.5 extracted the text. Super 3.5 VL corrected it from an enlarged crop of each block. It can still contain recognition errors.
Page headers, footers, and page numbers are removed.

\subsection*{Untitled continuation (E0181)}
\textit{Continues an entry from an earlier page.}

escape or advance, he was compelled to sign a capitulation by which the whole army laid down their arms, and were differenced into different quarters of cantonement. By this remarkable capitulation, which was called the capitulation of Closter Seven, Hanover was obliged to submit quietly to the French, who were now determined to turn their arms against the king of Prussia.

Soon after this capitulation, both sides began to complain that the treaty was not strictly observed. The Hanoverians exclaimed against the rapacity of the French general, and the brutality of his soldiers. The French retorted the charge against them; accused them of insolence and infraction; and, being sensible of their own superiority, refused to bind them strictly to their terms of agreement. The Hanoverians only wished for a pretence to take arms, and a general to head them. Neither were long wanting. The oppressions of the tax-gatherers whom the French had appointed, were considered as so severe, that the army rose to vindicate the freedom of their country, while Ferdinand, prince of Brunswick, put himself at their head. As soon as this was known in Britain, large supplies were granted both for the service of the king of Prussia, and to enable the Hanoverian army to act vigorously in conjunction with him. A small body of British forces was sent over to join prince Ferdinand under the duke of Marlborough. After some inconsiderable successes at Crevelt, the duke of Marlborough dying, the command of the British forces devolved on lord George Sackville. A misunderstanding arose between him and prince Ferdinand, which appeared at the battle of Minden that was fought shortly after. Lord George pretended that he did not understand the orders sent him by the prince, and of consequence did not obey them. The allies gained the victory, which would have been more decisive had the British commander obeyed his orders. He was soon after recalled, tried by a court-martial, found guilty of disobedience, and declared incapable of serving in any military command for the future.

After this victory it was imagined, that one reinforcement more of British troops would terminate the war in savour of the allies; and that reinforcement was quickly sent. The British army in Germany was augmented to upwards of 30,000 men, and sanguine hopes of conquest were generally entertained. These hopes, however, were soon found to be ill founded. The allies were defeated at Corbach; but retrieved their honour at Exdorf. A victory at Warburgh followed shortly after, and another at Zierenberg; but then they suffered a defeat at Compen; after which, both sides retired into winter-quarters.

On the 25th of October 1760, happened the death of king George II. He had risen at his usual hour, and observed to his attendants, that as the weather was fine, he would take a walk into the gardens of Kensington, where he then refreshed. In a few minutes after his return, being left alone, he was heard to fall down upon the floor. The noise of this bringing his attendants into the room, they listed him into bed; where he desired with a faint voice, that the princess Amelia might be sent for: but before she could reach the apartment, he expired, in the 77th year of his age, and 33d of his reign. An attempt was made to bleed him,


\section*{Source record}
\begin{description}
\item[Scan ID] \texttt{britannica:309f765b1ecc83e0a7f7762e44b4a927a7555443:train:2108}
\item[Image file] \texttt{page-2108-6874128a519b.jpg}
\item[Image SHA-256] \texttt{6874128a519b8cf9cd9de001ba48c06556b084267e667c07a65ea596b6b51576}
\item[Dataset] NationalLibraryOfScotland/encyclopaedia\_britannica\_illustrated at revision \texttt{309f765b1ecc83e0a7f7762e44b4a927a7555443} (CC0-1.0)
\item[Extraction] model \texttt{nvidia/nemotron-3.5-parse}, raw reply \texttt{extract/raw/page-2108-6874128a519b.json}
\item[Raw reply SHA-256] \texttt{2a03288f9ee4f3341c3f0ce34e3c33bcb383d7be7326b17e3628b368e7983bc0}
\item[Correction] model \texttt{nvidia/Nemotron-3.5-Super-VL}, block crops, raw replies \texttt{correct/raw/crop\_nothink/page-2108-6874128a519b-b*.json}; super\_vl: 4
\item[Running header] B R I  [ 645 J B R F
\end{description}
\end{document}
